\section{Implementación}
\subsection{Explicación del solver}
lp\_solver es un solver y biblioteca que resuelve problemas de programación lineal y programación entera mixta utilizando los métodos simplex y branch and bound.
\newline Tras recibir un modelo, incluyendo variables de decisión, función objetivo y restricciones, de parte del usuario, que puede ser desde archivos de texto, archivos MPS o directamente mediante llamadas a su propia API, el modelo puede, en primer lugar, aplicar un preprocesamiento llamado \"Presolve\" que simplifica el modelo, eliminando variables y restricciones que no aportan al problema y sustituyendo variables con valores fijos por constantes, reduciendo la complejidad a nivel computacional del problema.
\newline En segundo lugar, lp\_solve puede usar métodos de escalado (scaling) para cuando el modelo involucra coeficientes con magnitudes muy pequeñas y/o muy grandes. Reduce temporalmente las magnitudes a valores idealmente cercanos 1 para mejorar la estabilidad numérica y reducir errores de redondeo.
\newline Luego, comienza con la solución del problema utilizando una variante del método simplex denominada método simplex revisado, que utiliza una representación matricial y solo calcula la información necesaria para cada iteración.
\newline Si tras utilizar el método anterior se llega a una solución óptima continua cuando existen restricciones de valores enteros en el modelo, el solver pasa a utlizar branch and bound como se conoce, diviendo el problema en subproblemas que abarca cada uno una restricción adicional que fuerza a las variables a tomar valores enteros, resolviendo cada subproblema por separado y descartando la rama a medida que el subproblema no genere una solución más eficiente.
\newline Para finalizar, lp\_solve entrega una serie de resultados sobre el problema, incluyendo si se llegó a una solución factible o no, el valor de la función objetivo, los valores asignados a las variables de decisión, el cumplimiento de las restricciones y el tiempo de ejecución.
